\begin{table}[t]
\centering
\caption{\textbf{Why the Hawkes popularity prior runs away on KuaiRand but not
Micro-Video or MovieLens} (identical SASRec+Ours architecture/framework, fit
independently on each dataset). Two natural hypotheses -- denser per-item interactions,
or more skewed raw popularity -- are each contradicted by at least one other dataset
(bold). What tracks the outcome instead is a short absolute time span combined with high
local density concentrated in a few items, which forces a very fast shared decay rate
$\beta$ and, combined with a fixed 50-event-per-item history cap, leaves the excitation
channel unusable for the vast majority of a popular item's own history.}
\label{tab:kuairand_dataset_comparison}
\begin{tabular}{l ccc}
\toprule
 & KuaiRand & Micro-Video & MovieLens \\
\midrule
catalog size ($|\mathcal{V}|$) & 7{,}076 & 44{,}503 & 3{,}533 \\
mean interactions / item & 70.1 & 4.3 & $\mathbf{132.7}$ \\
raw count max/median ratio & 323 & $\mathbf{404}$ & 60 \\
\midrule
dataset time span (days) & $\mathbf{29.5}$ & 27.8 & 1{,}038.8 \\
top item's raw count & 5{,}704 & 404 & 2{,}508 \\
$\Rightarrow$ top item's interactions/day & $\mathbf{\sim}193$ & $\sim$15 & $\sim$2.5 \\
median same-item re-interaction gap & $\mathbf{0.35}$h & 1.99h & 6.78h \\
learned $\beta$ (decay/day) & $\mathbf{18.04}$ & 6.22 & 0.107 \\
$\Rightarrow$ excitation half-life & $\mathbf{0.92}$h & 2.67h & 6.48\emph{d} \\
usable excitation window for top item\tnote{*} & $\mathbf{1.1}$d & 0.4d & 684.7d \\
\midrule
learned $\mu$ max/median ratio & $\mathbf{5.2\times10^6}$ & 1{,}727 & 2{,}673 \\
top-10 items' share of total $\mu$ mass & $\mathbf{87.3\%}$ & 7.1\% & 28.3\% \\
Tail Recall@10, backbone $\to$ Ours ($\eta{=}0.5$) & $\mathbf{0.0268\to0.0036}$ & $0.0112\to0.1449$ & $0.0365\to0.0451$ \\
\bottomrule
\end{tabular}
\begin{tablenotes}
\item[*] Days spanned by the last 50 raw event timestamps retained per item
(`time\_dict\_to\_array`, `time\_len=50`), for the item with the single largest learned
$\mu_v$ in each dataset.
\end{tablenotes}
\end{table}
